\newcommand{\figuraCilindroArquimedes}{%
\begin{figure}[H]
    \centering
    \begin{tikzpicture}[>=Latex,font=\small,line cap=round,line join=round]
        \draw[very thick] (1.0,5.0)--(1.0,0.4)--(10.8,0.4)--(10.8,5.0);
        \fill[blue!12] (1.08,0.48) rectangle (10.72,4.35);
        \draw[very thick,blue!60!black] (1.0,4.35)--(10.8,4.35);

        \fill[orange!25] (4.4,1.45) rectangle (7.0,3.65);
        \draw[very thick,orange!80!black] (4.4,1.45) rectangle (7.0,3.65);
        \draw[densely dashed] (4.4,3.65)--(7.0,3.65);
        \draw[densely dashed] (4.4,1.45)--(7.0,1.45);

        \draw[->,ultra thick,red!75!black] (5.7,3.68)--(5.7,2.82)
            node[midway,right] {$p_1A$};
        \draw[->,ultra thick,red!75!black] (5.7,1.42)--(5.7,2.28)
            node[midway,right] {$p_2A$};
        \node[anchor=east] at (4.25,3.65) {cara superior, $p_1$};
        \node[anchor=east] at (4.25,1.45) {cara inferior, $p_2$};

        \draw[<->,thick] (7.45,1.45)--(7.45,3.65)
            node[midway,right] {$L$};
        \draw[<->,thick] (4.4,1.05)--(7.0,1.05)
            node[midway,below] {área $A$};

        \node[align=left] at (9.15,2.65)
            {$\displaystyle F_B=(p_2-p_1)A$\\[2mm]
             $\displaystyle =\rho_f gLA$\\[2mm]
             $\displaystyle =\rho_f gV$};
    \end{tikzpicture}
    \caption{Derivación del empuje mediante un elemento cilíndrico. Como la
    presión aumenta con la profundidad, la fuerza sobre la cara inferior
    supera a la fuerza sobre la cara superior.}
    \label{fig:cilindro-arquimedes}
\end{figure}%
}